
\documentclass{article}
\usepackage{colortbl}
\usepackage{makecell}
\usepackage{multirow}
\usepackage{supertabular}

\begin{document}

\newcounter{utterance}

\centering \large Interaction Transcript for game `cladder', experiment `full\_v1.5\_default', episode 9980 with Qwen3.5{-}4B.
\vspace{24pt}

{ \footnotesize  \setcounter{utterance}{1}
\setlength{\tabcolsep}{0pt}
\begin{supertabular}{c@{$\;$}|p{.15\linewidth}@{}p{.15\linewidth}p{.15\linewidth}p{.15\linewidth}p{.15\linewidth}p{.15\linewidth}}
    \# & \multicolumn{2}{c}{Player} && \multicolumn{2}{c}{Game Master} \\
    \hline

    \theutterance \stepcounter{utterance}  
    & & & \multicolumn{4}{p{0.6\linewidth}}{
        \cellcolor[rgb]{0.9,0.9,0.9}{
            \makecell[{{p{\linewidth}}}]{
                \texttt{\tiny{[P1$\langle$GM]}}
                \texttt{You are an expert in causal inference. The following question is not a typical commonsense query, but rather a meticulously designed question created by a professor specializing in causal inference, intended to assess the students' mastery of the course content.} \\
\texttt{Imagine a self{-}contained, hypothetical world with only the following conditions, and without any unmentioned factors or causal relationships: Hwax has a direct effect on pexu and rukz. Pexu has a direct effect on rukz. We know that hwax causes pexu. hwax and pexu causes rukz. We observed an individual is not hwax. Would an individual is not rukz if not pexu instead of pexu?} \\
\texttt{Start your answer with 'Yes' or 'No', followed by additional reasoning or evidence to support your explanation.} \\
            }
        }
    }
     \\ \\

    \theutterance \stepcounter{utterance}  
    & \multicolumn{4}{p{0.6\linewidth}}{
        \cellcolor[rgb]{0.561,0.737,0.561}{
            \makecell[{{p{\linewidth}}}]{
                \texttt{\tiny{[P1$\rangle$GM]}}
                \texttt{No} \\
\\ 
\texttt{**Reasoning:**} \\
\\ 
\texttt{To determine the answer, we must construct the causal structure based strictly on the provided conditions and apply the rules of counterfactual reasoning (specifically the "do{-}calculus" or structural causal models).} \\
\\ 
\texttt{1.  **Construct the Causal Graph:**} \\
\texttt{    *   \$Hwax \textbackslash{}rightarrow Pexu\$ (Given: "hwax causes pexu").} \\
\texttt{    *   \$Hwax \textbackslash{}rightarrow Rukz\$ (Given: "Hwax has a direct effect on rukz").} \\
\texttt{    *   \$Pexu \textbackslash{}rightarrow Rukz\$ (Given: "Pexu has a direct effect on rukz").} \\
\texttt{    *   *Note:* The statement "hwax and pexu causes rukz" is consistent with the graph where both are parents of \$Rukz\$.} \\
\\ 
\texttt{2.  **Analyze the Counterfactual Query:**} \\
\texttt{    The question asks: "Would an individual be not \$Rukz\$ if not \$Pexu\$ instead of \$Pexu\$?"} \\
\texttt{    *   **Context:** We observed an individual who is **not** \$Hwax\$ (\$Hwax = 0\$).} \\
\texttt{    *   **Intervention:** We intervene to set \$Pexu = 0\$ ("not \$Pexu\$").} \\
\texttt{    *   **Target:** Determine the outcome of \$Rukz\$ under this intervention.} \\
\\ 
\texttt{3.  **Evaluate the Mechanism:**} \\
\texttt{    *   In this world, \$Rukz\$ is determined by its parents: \$Hwax\$ and \$Pexu\$.} \\
\texttt{    *   Since we are in a hypothetical world with "only the following conditions," we assume a deterministic relationship where \$Rukz\$ is a function of its causes (\$Rukz = f(Hwax, Pexu)\$).} \\
\texttt{    *   We know \$Hwax = 0\$.} \\
\texttt{    *   We are intervening to set \$Pexu = 0\$.} \\
\texttt{    *   Therefore, the inputs to the function determining \$Rukz\$ are \$(0, 0)\$.} \\
\\ 
\texttt{4.  **Logical Deduction:**} \\
\texttt{    *   If \$Hwax\$ causes \$Rukz\$ and \$Pexu\$ causes \$Rukz\$, then \$Rukz\$ exists only if at least one of its causes is active (assuming a standard "OR" logic for causal sufficiency in such binary hypothetical worlds, or simply that the absence of causes leads to the absence of the effect).} \\
\texttt{    *   If \$Hwax\$ is absent and \$Pexu\$ is absent, there are no active causal forces driving \$Rukz\$.} \\
\texttt{    *   Consequently, \$Rukz\$ will not occur.} \\
\\ 
\texttt{5.  **Addressing the "Instead of" phrasing:**} \\
\texttt{    The phrase "} \\
            }
        }
    }
    & & \\ \\

    \theutterance \stepcounter{utterance}  
    & & & \multicolumn{2}{p{0.3\linewidth}}{
        \cellcolor[rgb]{0.9,0.9,0.9}{
            \makecell[{{p{\linewidth}}}]{
                \texttt{\tiny{[GM$|$GM]}}
                \texttt{No} \\
            }
        }
    }
    & & \\ \\

    \theutterance \stepcounter{utterance}  
    & & & \multicolumn{2}{p{0.3\linewidth}}{
        \cellcolor[rgb]{0.9,0.9,0.9}{
            \makecell[{{p{\linewidth}}}]{
                \texttt{\tiny{[GM$|$GM]}}
                \texttt{yes} \\
            }
        }
    }
    & & \\ \\

    \theutterance \stepcounter{utterance}  
    & & & \multicolumn{2}{p{0.3\linewidth}}{
        \cellcolor[rgb]{0.9,0.9,0.9}{
            \makecell[{{p{\linewidth}}}]{
                \texttt{\tiny{[GM$|$GM]}}
                \texttt{game\_result = LOSE} \\
            }
        }
    }
    & & \\ \\

\end{supertabular}
}

\end{document}
